
\documentclass{article} % For LaTeX2e
\usepackage{icais2025_conference,times}

% Optional math commands from https://github.com/goodfeli/dlbook_notation.
\input{math_commands.tex}

\usepackage{hyperref}
\hypersetup{
    colorlinks=true,
    linkcolor=blue!70!black,
    citecolor=green!60!black,
    urlcolor=red!60!black
}
\usepackage{url}

\documentclass{article}
\usepackage{amsmath}
\usepackage{amssymb}
\usepackage{array}
\usepackage{algorithm}
\usepackage{algorithmic}
\usepackage{algorithmicx}
\usepackage{booktabs}
\usepackage{colortbl}
\usepackage{color}
\usepackage{enumitem}
\usepackage{fontawesome5}
\usepackage{float}
\usepackage{graphicx}
\usepackage{hyperref}
\usepackage{listings}
\usepackage{makecell}
\usepackage{multicol}
\usepackage{multirow}
\usepackage{pgffor}
\usepackage{pifont}
\newtheorem{definition}{Definition}
\usepackage{soul}
\usepackage{sidecap}
\usepackage{subcaption}
\usepackage{titletoc}
\usepackage[symbol]{footmisc}
\usepackage{url}
\usepackage{wrapfig}
\usepackage{xcolor}
\usepackage{xspace}
% \usepackage[numbers]{natbib}  
\usepackage[utf8]{inputenc}
\usepackage{amsmath, amsfonts, amssymb}
% \DeclareMathOperator*{\argmax}{arg\,max}
\usepackage{graphicx}
\usepackage{hyperref}

\newenvironment{myquotation}
  {\begin{quotation}}  % begin code
  {\end{quotation}} 

\title{Economic Implications of Language Models and Copyright Law}

\begin{document}
\maketitle
\begin{abstract}
How will language models (LMs) affect future economic progress?
Inspired by the \emph{Lever of Riches} by \citet{mokyr1992lever}, we argue that the institutions governing LM content generation and usage patterns are critical to answering this question.
We content that, because LM creators have a strong incentive to collect, train, and deploy intellectual property protection, the all-you-can-consume access to knowledge and creativity they enable has led to rapid acceptance and widespread use, which in turn results in smaller, low-skill employment creation but increased output and greater overall welfare.
We provide a theoretical and analytical framework explaining this phenomenon and point to its long-term consequences using empirical evidence.
\end{abstract}

\section{Introduction}


\label{Intro}

With the advent of the fourth industrial revolution enabled by language models (LMs) like ChatGPT, there has been an increased interest in understanding their economic impact~\citep{Howard2019artificial, autor2015there, hui2024short, felten2023will, eloundou2023gpts, paslar2023role}.
As \citet{boiko2023emergent} suggest, the fourth industrial revolution---enabled by LMs---is different from previous revolutions in its ability to replace creativity and cognitive skills beyond what was previously possible; they argue that this makes the study of the economic impact of LMs both more challenging and more important to study.
The \emph{Lever of Riches} by \citet{mokyr1992lever} argues that technological innovation is a key driver of economic growth and development, stating that improvements in institutions and institutions that protect intellectual property (IP) are crucial levers for economic development.
Following the \emph{Lever of Riches}, we argue that the institutions governing LM content generation and usage patterns play a critical role in how they impact the economy.
Our research indicates that the rapid development and adoption of LMs are likely to lead to smaller, low-skill employment creation but increased economic progress and overall welfare. 

Specifically, we argue that the strong intellectual property of LMs, in combination with their ability to produce convincing content without obvious provenance, will lead to a rapid acceptance and use of LMs, resulting in economic progress.
At the same time, we argue that the weak intellectual property protection of the historical creators of content (such as writers and musicians) and the rapid automated generation and production of low-quality content, of such works like college essays~\citep{thompson2022chatgpt} and artworks~\citep{hullman2023artificial} are likely to result in slower technology adoption and development.
This dynamic is likely to create new, low-skill job opportunities in the digital economy, while shrinking more traditional creative and knowledge-intensive industries.
It will lead to greater output and consumer welfare, but also increased wealth inequality and higher costs for access to education and labor protection.
Lastly, we argue that changes in laws and regulations will have a significant impact on the speed and impact of LMs on the economy. 
The rapid development and adoption of LMs present governments with an important role to play in protecting the rights of historical creators and in ensuring that the economic impacts of LMs are fair and equitable. 
The findings of this paper can help policy makers and tech corporations adapt institutions to cope with the challenges presented by LMs. 
To summarize, our contributions are as follows:

    (i) We identify intellectual property laws as the main institutions influencing the speed and direction of LM adoption and its potential economic impact.
    (ii) We provide a theoretical framework explaining how this impact might manifest over time.
    (iii) We provide evidence of this impact via case studies and trends in a wide range of domains.



\section{The \emph{Lever of Riches}}

\label{Prelim}

\emph{Leverage}, by \citet{mokyr1992lever}, provides a compelling argument in favor of the role of technology in economic development.
The \emph{Lever of Riches} is an economic theory proposing that technology---as the key driver of economic growth and development.
According to the \emph{Lever of Riches}, the introduction of a new and ever-improving technology results in economic progress and development.
This occurs through:

    (i) \emph{The creation of new opportunities}~\citep{mokyr1992lever, autor2003skill}:
    New technologies lead to the emergence of new opportunities for individuals to engage in profitable business activities.
    For example, the invention and spread of the World Wide Web fostered the growth of web developers and web designers.
    (ii) \emph{The reallocation of factors of production}~\citep{autor2003skill, autor2015there, restrepo2023automation}:
    New technologies necessitate a reallocation of resources to more productive activities.
    (iii) \emph{The reallocation of labor to more productive tasks:} New technologies allow individuals to engage in more productive tasks. The result is an increase in economic output.


The creation of new opportunities results in the emergence of new industries and jobs, following the \emph{creative destruction} principle as outlined in \citet{aghion1990model, schumpeter2013capitalism, autor2003skill, autor2015there}:   \emph{Competitive capitalism, in fact, does allow a society which is neither progressive nor stationary. It may be progressive for an unlimited time if the efficiency units in the production of marketable commodities by means of available factors of production are never allowed to coincide with the actual units employed.}\footnote{\citet{aghion1990model} delineated a model of growth driven by continuous innovation, emphasizing the role of the destruction of comparative advantage as a catalyst for growth.}

\citet{smith2002inquiry} makes a similar argument on the role of innovation, stating:   \emph{The ability to achieve this high degree of regular and equal admittance (to the benefits of economic discovery) is indeed the principal cause of the productive and enriching effects which accompany a division of labour.}\footnote{\citet{smith2002inquiry} emphasizes the role of labor division and equal access in fostering productive economic outcomes.}

However, there are significant political economy challenges regarding such developments, particularly because of the destruction of comparative advantages held by individuals.
The \emph{Gerschenk effect}~\citep{autor2003skill} argues that technological advances cause temporary job loss for specific groups. 
\citet{Friedman1957ConsistencyOT} argues that new technologies lead to skill complementing and destruction, where routine and mid-level skills are most affected.
This is further supported by the Law of Comparative Advantages~\citep{smith2002inquiry}:     \emph{The greater the variety of capabilities that one can draw on from many individuals, the greater the likelihood of complementary skills and the emergence of labor specialization.}\footnote{\citet{smith2002inquiry} highlights the importance of a diverse and specialized workforce for productivity gains.}

The \emph{prisoner's dilemma}~\citep{poundstone2011prisoner} and \emph{Coase theorem}~\citep{coase2013problem} argue that individuals may not have a collective incentive to cooperate and establish institutions that regularize their interactions and allow for sharing and leveraging of comparative advantages.
This is because individuals often lack collective action capacity and may not support institutions that could enable greater collaboration and information sharing~\citep{buchanan1965calculus, olson2012logic, ostrom1990governing}.
Thus, government intervention is often necessary to facilitate such collaborations.

The arguments of the \emph{Lever of Riches} and the \emph{Gerschenk effect} and Laws of Comparative Advantages~\citep{smith2002inquiry} are built upon the microdynamics of technological changes on individuals and their productivity.
A key---though controversial~\citep{landymore2023sam, finance2024ai}---{argument} in support of these theories is  the development of ever-improving technologies necessarily leads to increased output, which implies a general economic improvement.
However, there are significant arguments against this \emph{harper collapse} in the \emph{Lever of Riches} from both a \emph{long-run} and \emph{short-run} perspective~\citep{acemoglu2023power, swoboda2025examining}.
\citet{acemoglu2023power} argues that technological diffusion can exhibit diminishing marginal returns in the \emph{long-run}, potentially leading to decreased output.
Others have noted long-term (\emph{endogenous}) AI existential risks (\emph{EARs})~\citep{bender2021dangers, boiko2023emergent, lu2024ai, Kulveit2025GradualDS, hendrycks2023overview}.
From a short-run perspective, even if technological improvements continue, the impact of increasing labor market tightness (due to a low inflation labor market and decreasing labor force growth) and skill bias technological improvements may slow down the velocity of economic progress due to a decline in labor \emph{utilization rate}~\citep{acemoglu2012why}.
\citet{acemoglu2012why} document this decline over the past two decades in advanced economies.


\section{Formalizing Intelligent Economic Agents}

\label{RfC}

The aforementioned arguments about what economic impacts technological advances may have rely on key assumptions about individuals' behaviors.
Often, theories and models are built on the principle of \emph{earning maximization} of \emph{self-interested agents}.
This implies that individuals are:    \emph{interested in their own profit; they are not eager to communicate their own knowledge to others, but only to use the knowledge of others for their own purposes,} \emph{since such communication is costly.}\footnote{\citet{smith2002inquiry}\emph{Smitherus:} an individual's inclination to pursue self-interest, employing knowledge for personal gain while reluctant to share their own knowledge.}

\emph{Earning maximization} implies that individuals are \emph{utilitarian} and therefore interested in the \emph{marginal revenue product} of their labor, i.e., the difference between revenue (the amount users are willing to pay for a job) and the cost of labor (the amount one has to pay factors of production, e.g., wages).

To see this, consider a farmer who owns a field and wants to produce crops from it.
\begin{definition}[{\bf Marginal Revenue Product}]\label{def:mrp}
    Given a production function, $y=\mathcal{F}(e)$ determining the output $y$, where $e$ is a vector of factors of production (e.g., labor, capital), the \textit{marginal revenue product} for each factor of production $F_i$ is defined as:
    \begin{equation*}
        MRP_i(y) := \textup{ Marginal} \; \textup{Contribution} \left(F_i \;  \lvert \; \mathcal{F}, \right)  = \left(\frac{\partial \mathcal{F}(e) }{\partial F_i}\right)
    \end{equation*}
\end{definition}

Let's say that the farmer owns 180 acres of land, and farming 1 acre yields 2 units of crops, while the cost of labor (wages, equipment, etc.) is \$10 per acre.
The farmer can, thus, produce $360=180\cdot 2$ units of crop, while incurring a cost of $3,600=10\cdot 180$ and thus make a profit of $2,100=3,600-3,600$ dollars by farming the owned land.
However, suppose another farmer owns 185 acres of land.
The second farmer can either farm 5 acres, or farm 180 acres to farm 180 acres and sell.
Farming 5 acres yields 18 units of crop, with a profit of $5=10\cdot 5$, while farming 180 acres yields $360=180\cdot 2$, with a profit of $1,820=10\cdot 180-3,600$.
Thus, farming 180 acres is profitable, but farming 5 acres is not.
Based on this, a farmer with 185 acres of land would farm 180 acres, and not 185 acres, because the marginal revenue product for land is lower.
This also explains why farming 5 acres yields less revenue and profit than farming 180 acres, despite using 35 fewer acres of land.
This also suggests that, if the price of factors of production for farming---such as labor, equipment, machinery, and land---increases by 5\% and crop production remains the same, farming 180 acres might not be profitable, which implies farmers would reduce farm size, leading to decreased output and increased wealth inequality due to higher factors of production costs.

The \emph{marginal revenue product} can be expanded to incorporate other factors of production. 
If the price of factors of production increases, farmers will have to find ways to produce more units given the same amount of resources.
To see this, let's look at the farmer with 185 acres of land. 
If the price of a crop decreases by 1\%, or each crop unit is now worth 12 instead of 13, and the price of factors of production (labor, equipment, machinery, and land) increases by 10\%, while the cost of factors of production for producing 180 acres (the land owned by the first farmer) becomes 18,000, the farmer still makes a profit of 8,400, the same as the first one, by farming the entire 185 acres of land.
This also suggests that farmers will have to innovate and find ways to farm these acres cheaper if they are to make a profit.
Now, let's consider the farmer with 185 acres of land.
If the price-to-earnings ratio for a crop decreases by 1\%, each crop unit is now worth only 12, and factors of production for farming become more expensive, the farmer can no longer farm the entire 185 acres of land and make a profit.
Instead, she would have to find a way to farm a smaller amount of land while still making a profit.
This also suggests that farmers may need to innovate to reduce labor costs or reduce the amount of land needed to farm a unit of crop.

\subsection{Defining LM Productivity}

Based on this, we can now define what we mean by \emph{productivity enabled by LMs}.
Consider a user $u$ with task specification $P_u$, such as "Write an art essay on Hamlet," where $P_u \in P$, the set of all possible task specifications.
Let's denote $Cost_u(x)$ as the (marginal) cost of producing some output $x$ for user $u$ using factors of production $F$.
For example, these factors of production may include the amount of labor, equipment, machinery, software licenses, electricity, and other costs (such as the OpenAI API cost of using ChatGPT or the electricity needed for running your computer at home for the "Lever of Riches" analysis).
We denote the utility function of user $u$ as $v_u: P \times \mathcal{N} \to \mathbb{R}$, which is a combination of a take-or-leave function $w_u: \mathcal{N} \to \mathbb{R}$, where $w$ is the write functionality of an LM, that takes in prompts as input and produces outputs, specifying what users enjoy from the LM and a try-or-leave function $l_u: \mathcal{N} \to \mathbb{R}$, where $l$ is the edit or refinement functionality of an LM, that takes the outputs produced by $w$ and further modifies them to better fit user $u$'s goal:    \emph{such functions may continuously refine the response provided by ChatGPT to better align with user $u$'s objectives or preferences.}\footnote{\citet{salvi2024conversational} introduces a method to fine-tune LMs for specific writing tasks based on user feedback and preferences.}

The output produced by taking the write function, $w_u(P_u)$, of user $u$ with task specification $P_u$ is denoted as $x_u \in \mathcal{N}$, where $\mathcal{N}$ is the output space of $w$.
For example, $w_u$ may produce an art essay on Hamlet based on $P_u$.
The goal of user $u$ is to edit (or modify) $x_u$ to best achieve their task goal $P_u$.
The costs of doing so, denoted $Cost_u(x_u)$, may include the marginal cost of labor (typing or speaking) for editing.
The quantity $x^*_u \in \mathcal{N}$ denotes the output (e.g., essay) that user $u$ is aiming for.
Given the definitions above, the productivity of user $u$ is defined as:
\begin{equation}
    Pro_d(u, P_u) = \mathbb{E}_{x \sim D (x^*_u)} \left[ v_u \left( x \right) \right] - Cost_u \left( x \right),
\end{equation}
where $D$ is the distribution of outputs that user $u$ aims for and $pro_d: (u, P_u) \to \mathbb{R}$ is the expected utility achieved by user $u$ with task specification $P_u$, and the quantity $pro_d(u, P_u)$ goes back to the marginal revenue product defined in Definition~\ref{def:mrp}.\footnote{To see this, note that $pro_d(u, P_u)$ is the difference between the expected utility, under some distribution of ideal outcomes, and the cost. Multiplying by the number of users for task $P_u$, this will determine the total output yielded by.}

We define $u$'s \textit{productivity gain} over time as $Pro_d(u, P_u^{(t)}) - Pro_{d-1}(u, P_u^{(t-1)})$, where $P_u^{(t)}$ denotes the task that user $u$ is working on or being committed to in time $t$.
For example, $Pro_d(u, P_u^{(t+1)}) - Pro_{d-1}(u , P_u^{(t+1)}) > 0$ means that user $u$ is more productive at time $t+1$ than time $t$ for task $P_u^{(t+1)}$.
For $t=0$, we have the initial \emph{marginal revenue product} by the set of factors of production $F_0$.
For $t>0$, $F_t$ might be more \emph{costly} than $F_0$ (e.g., OpenAI increases prices for their API) or less \emph{costly} (e.g., a user can do it at home for free and derive a higher utility from it (the direct utility $v_u$)); $F_t$ might be more or less \emph{effective} (i.e., better at producing value) than $F_0$, or not.
Thus, the increase in productivity gains over time might be slow or fast, depending on how these factors of production interact with each other and the utility function of the user.
Given $F_t$ and $w_u(P_u)$, user $u$ aims to achieve their goal by editing $x_u$ to be as close as possible to their aim:
\begin{equation}
    x^*_u = \argmax_{x \in \mathcal{N}} u \left( x, P_u, F_t \right)
    \label{Eq:x-u-star}
\end{equation}
where $u$ is a (task- and user-specific) utility function that depends on the write functionality $w$ of the LM, the factors of production $F_t$, and the task specification $P_u$.
Each user $u$ maximizes their expected utility, which depends on factors of production $F_t$ and task specification $P_u$, by choosing the best edit $x^*_u \in \mathcal{N}$ from the output produced by the take-or-leave function $w$, defined in Equation~\ref{Eq:x-u-star}:
\begin{equation}
    u\left( x, P_u, F_t \right) = v_u \left( x \right) - Cost_u \left( x \right)
\end{equation}


\section{Economic Implications of Intellectual Property Laws}

\label{IPlaws}

The analysis presented in this section is predicated on the assumption that users enjoy using ChatGPT and other LMs more than they enjoy writing or creating art.
We justify this assumption by stating that using LMs is \emph{the easiest pathway to achieving one's goals}, especially given the low opportunity cost of labor for low-skilled employees~\citep{autor2015there}.
The productivity gains achieved by using LMs are expected to be significant, as they enable users to bypass the difficulty of writing or creating art~\citep{langley00}.
However, as we will see over time, users will become \emph{much more} productive and enjoy \emph{greater} utility by using LMs and successfully achieving their goals.

\subsection{Analyzing Current Intellectual Property Laws}

The rapid adoption and use of LMs is only possible if    \emph{new ideas do not require the permission of others to be used as part of the process of creating novel ideas.}\footnote{\citet{von2006democratizing} highlights the importance of the ability to use and incorporate novel ideas without permission in the creation of innovation.} Intellectual property laws, in particular current copyright laws, enable the free use of ideas and techniques offered by LMs~\citep{ginsburg2025humanist, restrepo2023automation}:    \emph{ Copyright law does not protect theideas, procedures, or algorithms embodied in or expressed in \textit{(a)} processes, systems, machines, mechanisms, models, . . . , computer programs, computer components, computer components, . . . ,  . . . LMs . . . are not protected.}\footnote{\emph{Because ideas, algorithms, and computer programs are not protected by the copyright, the use of a \emph{specific expression} of an idea, algorithm, or computer program (such as a specific computer program and its documentation or a specific implementation of an idea or algorithm) is not prohibited.}} The strong intellectual property protection of LM creators and users enables the rapid development, adoption, and usage of LMs:

\paragraph{Development:}  The rapid development and adoption of LMs are only possible due to the availability of large amounts of data and computing capacity.
The development of LMs requires extensive data collection, cleaning, and annotation~\citep{neurips1, dodge-etal-2021-documenting}.
For example, OpenAI's ChatGPT was trained on a combination of books, software source code (thus patent-protected), and other text and code from the internet~\citep{langley00}.
The creators of LMs thus have a strong incentive to protect their data and training code from unauthorized use and piracy.
The creators of LMs also have a strong incentive to protect the codebase itself~\citep{henderson2023foundation, lemley2020fair}.
Thus, the developers and owners of datasets and models have strong incentives to \emph{protect} their use in model development.
Copyright law gives the right to recoup the cost of writing, recording, producing, collecting, or creating code or content, according to market prices.
This is important because the more beneficial this is to an LM creator, the greater their expected return will be if they are allowed to own the copyright to their creations, just like a software developer or music producer.
Rather, LMs allow their creators to own the copyright title.

This is key to the rapid development, adoption, and and development of LMs because:
    \emph{ Copyright protection is nontransferable: if you create a copyrighted work, only you can give anyone the permission to use it - not even if you sell your business to another publisher.}\footnote{\citet{carlini2023extracting}: \emph{this is why major corporations spend hundreds of millions on exclusive licensing contracts to keep their competitors from using their trained data.}
}

This is important because the more beneficial it is to an LM creator, the greater their expected returns will be if they are allowed to own the copyright to their creations, just like a software developer or music producer.
Rather, LMs allow their creators to own the copyright title.

Copyright protection is beneficial to LMs' developers and owners because:    \emph{this is key to the rapid development, adoption, and development of LMs because:}
    \emph{Copyright protection is not transferable: if you create a copyrighted work, no one else can give anyone permission to use it.} [C]opyright protection lasts for life plus 70 years and is thus a strong incentive for creators to own their copyrights. This is important because the more beneficial it is to an LM creator, the greater their expected returns will be if they are allowed to own the copyright to their creations, just like a software developer or music producer. R

Copyright protection facilitates the rapid development, adoption, and and development of LMs because:    \emph{This is key to the rapid development, adoption, and development of LMs because ...} \emph{Copyright protection is not transferable: if you create a copyrighted work, no one else can give anyone permission to use it.} [C]opyright protection lasts for life plus 70 years and is thus a strong incentive for creators to own their copyrights. This is important because the more beneficial it is to an LM creator, the greater their expected returns will be if they are allowed to own the copyright to their creations, just like a software developer or music producer. R
\emph{Without the ability to own the copyright to their creations, LMs would not have access to the necessary data and copyrighted works, making their development and use much more difficult, if not impossible.}

This is because collecting, cleaning, and annotating data and then training models is very expensive.
Training large LMs is a complex and resource-intensive process, often requiring large datasets, sophisticated hardware, and significant energy consumption.
For example, large LMs like ChatGPT require thousands of GPUs, advanced software, large training datasets, and substantial computing power.
The time and resources needed for training also mean that LMs tend to be expensive to create, maintain, and run.
The law also gives creators a strong incentive to protect their code and model weights~\citep{henderson2023foundation, lemley2020fair}.

Copyright laws facilitate the rapid development, adoption, and and development of LMs because:    \emph{This is key to the rapid development, adoption, and development of LMs because ...} \emph{Copyright protection is not transferable: if you create a copyrighted work, no one else can give anyone permission to use it.} [C]opyright protection lasts for life plus 70 years and is thus a strong incentive for creators to own their copyrights. This is important because the more beneficial it is to an LM creator, the greater their expected returns will be if they are allowed to own the copyright to their creations, just like a software developer or music producer. R \emph{Copyright laws facilitate low prices for users because ...} \emph{Without the ability to own the copyright to their creations, LMs would not have access to the necessary data and copyrighted works, making their development and use much more difficult, if not impossible.} \emph{Copyright laws facilitate low prices for users because ...}

Intellectual property laws---in conjunction with market forces---also enable the rapid development, adoption, and and development of LMs because: \emph{Copyright laws facilitate low prices for users because ...} \emph{The open use and availability of such technologies creates a competitive market for AI-powered products and services, thus controlling potential market power and mitigating potential monopoly rents.} By protecting intellectual property associated with LMs, government thus \emph{creates collective incentives} for rapid technological development, adoption, and development. \emph{Copyright laws also facilitate low prices for users because ...} \emph{By protecting intellectual property associated with LMs, government thus creates collective incentives} for rapid technological development, adoption, and development.


Intellectual property laws create collective incentives (\emph{cross-subagent correlations}) for rapid technological development, adoption, and development because:  \emph{Intellectual property laws create collective incentives because ...}  For an \emph{individual}, the necessary factors of production $F$ (computing capacity, electric power, software licenses, etc.) for producing an output $x$, following Definition~\ref{def:mrp}, one can see that one is interested in maximizing marginal revenue product (i) if it becomes more expensive to produce the same output (ii) if the same output yields a higher revenue or utility. Given (i) or (ii), individuals are interested in reducing the cost of production. One way of doing so is by \emph{using ideas and techniques developed by others}.

The previous paragraph is a restatement of the \emph{Leisure Class theory}.\footnote{an argument for intellectual property}
Intellectual property laws are important because:    \emph{If the laws of a capitalists economy were so changed that it became more difficult for individuals to use one another's productivity levels, then ... Then individuals would have to ``work harder'' to achieve the same output. Indeed, in an extreme case, they might have to do everything from ``scratch'' each time they performed any output, and thus there would be no inducement to produce anything else than the bare minimum necessary for consumption.}


Intellectual property laws are beneficial to LMs because:  \emph{Intellectual property laws are beneficial to LMs because ...} \emph{Intellectual property laws create collective incentives} because individuals are interested in reducing the cost of production. One way of doing so is by \emph{using ideas and techniques developed by others}. \emph{Intellectual property laws are also beneficial to LMs because ...} \emph{If the laws of a capitalists economy were so changed that it became more difficult to use one another's productivity levels, then ... Then individuals would have more difficulty using LMs, making LMs more expensive to use than before.} \emph{Thus, intellectual property laws are important because ...} - given (i) or (ii), above) individuals are interested in reducing the cost of production. One way of doing so is by \emph{using ideas and techniques developed by others}.


The fact that individuals are interested in reducing the cost of production is also highlighted by the fast adoption of OpenAI's GPT-4, despite the \$12--100\$ per 1K tokens opened a market for products like GPTZero, an AI detection company and software provider that offers an AI detection service or plugin that flags any text likely to have been produced by AI.\footnote{\citet{gptzero_technology}: \emph{GPTZero's AI detection technology flags text likely generated by AI using two key methods: Features Extraction and Feature Filtering. First, it \textit{extracts distinctive features from the text}. These features are created by the unique patterns and structures in the text. GPTZero's technology uses these features to create a ``feature vector'' of the text. \dots without access to the feature vector since we don't have direct access to the AI's output. .}}
\citet{hazramartahallucination} shows that signature verification can improve LMs' detection.
\citet{hansspotting} propose a method for detecting LMs by analyzing text and other signals (e.g., word frequency and distribution in the text).
\citet{argyle2023out} use LMs (GPT-4) to simulate human survey responses, making it difficult to detect and leading to false conclusions.
\citet{wang2024economic} propose a strategy to reduce detection.

The fast adoption of OpenAI's GPT-4, despite the \$12--100\$ per 1K tokens, has opened a market for products like GPTZero, an AI detection company and software provider that offers an AI detection service or plugin that flags any text likely to have been produced by AI.\footnote{\citet{gptzero_technology}: \emph{GPTZero's AI detection technology flags text likely generated by AI using two key methods: Features Extraction and Feature Filtering. First, it \textit{extracts distinctive features from the text}. These features are created by the unique patterns and structures in the text. GPTZero's technology uses these features to create a ``feature vector'' of the text. \dots without access to the feature vector since we don't have direct access to the AI's output. .}}
\citet{hansspotting} propose a method for detecting LMs by analyzing text and other signals (e.g., word frequency and distribution in the text).
\citet{argyle2023out} use LMs (GPT-4) to simulate human survey responses, making it difficult to detect and leading to false conclusions.
\citet{wang2024economic} propose a strategy to reduce detection.

\textbf{Takeaways:}

    (i) The rapid development, adoption, and development of LMs is only possible due to \emph{strong} IP protection, which enables the open use of data, code, and insights and \emph{low} IP protection, which enables the spread, usage, and monetization of LMs.
    (ii) The market for products like GPTZero indicates a demand for lower IP protection.
    (iii) Intellectual property laws create a collective incentive for the fast adoption and use of LMs because individuals are interested in reducing the cost of production and are thus interested in using the work of others.
    (iv) Government intervention in the form of intellectual property laws is crucial in creating collective incentives for rapid technological development, adoption, and development.



\subsection{Analyzing Proposed Intellectual Property Laws}

Tech firms and governments have proposed changes to current copyright laws related to AI to enable corporations to continue collecting, training, and spreading LMs.
Some works argue that IP laws should be strengthened for tech firms because unauthorized use and access to data, software licenses, and content hinder the development of models that are perceived as \emph{underperforming}.\footnote{ \emph{But in machine learning, there is no such thing as a protected “specialized purpose” technology; everything is on the internet, open to general use, and subject to normal copyright laws.}}

There are other works argue that the current laws regarding IP in AI should be changed, in particular that \emph{ideas} should also be protected by copyright or copyright laws that protect LM outputs should be relaxed.
By doing so, they argue that one can also allow for a more competitive market to emerge and drive innovation.
\textbf{Takeaways:}
    (i) The rapid development, adoption, and development of LMs is only possible due to \emph{strong} IP protection, which enables the open use of data, code, and insights and \emph{low} IP protection, which enables the spread, usage, and monetization of LMs.
    (ii) The market for products like GPTZero indicates a demand for lower IP protection.
    (iii) Intellectual property laws create a collective incentive for the fast adoption and use of LMs because individuals are interested in reducing the cost of production and are thus interested in using the work of others.
    (iv) Government intervention in the form of intellectual property laws is crucial in creating collective incentives for rapid technological development, adoption, and development.

\section{Economic Evidence: Documenting the Trade-offs}

\label{Evidence}

\subsection{Case Studies: Documenting the Trade-offs in Practice}

We begin by examining the impact of AI on a range of occupations, from creative industries (e.g., artists, musicians)---where AI has long been seen as a potential threat to jobs---to online platforms and medical professionals---areas where AI is transforming in ways that were previously unimaginable.
The evidence \emph{confirming} these trends is also emerging: LMs are rapidly becoming a fundamental part of everyday life, extending beyond human-created content (such as books, art, and music) to online platforms and services. 
In fact the use of GPT-4 API increased by $138\%$ and more than doubled the number of users from 2021 to 2023, with 2023 being the period during which we collected data for this paper. 
\paragraph{Online Platform Workers: The Trade-off \emph{Always} Occurs.}
We first focus on the widely discussed trade-off between the rapid adoption of LMs by highly paid workers in industries such as software and video programming, design, and business and management consulting, on the one hand, and the displacement of low-wage workers in industries such as content creation and customer assistance, on the other.
This trade-off is well-documented by the fast adoption of LMs in combination with the displaced human workers who are replaced by LMs.
Human workers are replaced by LMs especially in low-skill and knowledge-intensive industries.
The evidence for this can be seen in lawsuits and law firm representations, research studies and articles, and news articles and media coverage.
Often, the evidence comes in the form of lawsuits and law firm representations, as discussed below in Section.
For the first category, we find evidence that LMs are becoming a fundamental part of everyday work in a wide variety of fields, with a trend towards replacing human workers in low-skill, high-knowledge-intensive occupations.
For the second category, we compile evidence of the rapid adoption of LMs especially by high-wage and high-skill workers.
\paragraph{Online Platform Workers.}
Online platforms, such as Amazon Mechanical Turk\footnote{\url{https://www.mturk}}. While there is evidence that LMs lead to a reduction in online work participation among majority groups, there also appears to be an increase among in minority groups.
This suggests that platforms may rely on LMs to still get the work done, but they no longer need (or at least not as many) workers.
This also means that platform workers receive even lower pay and are more susceptible to significant pay cuts.
\paragraph{Medical and Health Workers.}
The use of AI in this sector is continually rising, with over 25\% of healthcare professionals worldwide believing that AI will increase their job security, with the potential for over 30\% of healthcare jobs to be replaced by AI.
However, only 11\%  believe that their pay will increase correspondingly. 
Others also finds that 43\% of healthcare professionals believe their employment conditions will worsen. 
Other work has also documented similar trends, with \citet{morley2020ethics} highlighting the potential negative impacts of AI on government-funded healthcare systems due to its potential to dehumanize healthcare, negate accountability, and undermine healthcare access for minorities.
Similarly, \citet{rony2024wonder} finds that medical professionals are particularly concerned about the impact of AI in their industries, with nurses being the most worried occupation in the working age population in the US. 
\paragraph{Writers, Journalists, Essayists.}
The adoption of LMs in this sector is most prominent, with over 50\% of college students with access to the internet reporting the use of LMs for essay writing~\citep{harvard2025chatgpt}, with only 9\% of students reporting that their professors prohibit or have warned against the of use LMs in essay writing~\citep{si2024can}.
At the same time, over 30\% of professors believe that the use of LMs results in the generation of weak, unoriginal arguments, with nearly 50\% of survey respondents believing that LMs negatively affect both the originality and quality of student work~\citep{porter2024ai}.
In the long run, this is expected to reduce (already falling) writing scores in essays, journalism, and other writing-related tasks~\citep{moon2024homogenizing}.
\paragraph{Programmers, Software Illustrators, Software Developers, Game Designers.}
LMs have a large potential market in programming, with 80\% of developers with 5 years of experience or more knowing, using, or recommending LMs, with this number increasing with the number of programming years~\citep{peng2023impact}.
This is not surprising, as LMs are especially well-suited to software design, management, and coding.
In fact, code and software source, which are key resources for developing LMs, are themselves mostly not protected by copyright~\citep{henderson2023foundation}.
Thus, as discussed in Section~\ref{IPlaws}, LM developers and owners have a strong incentive to protect the code and source for their databases, model weights, and models.

\subsection{Market Power: The Trade-off \emph{Always} Happens}

The concentration of corporations in industries like software, video, and board games is (expectedly) having a large impact on the fast adoption of LMs and consequently on labor markets.
The top four companies of the world for software (openAI, Google, Microsoft, and Cisco) video (Walt Disney, Sony, Comcast, and Amazon), and board games (Hasbro, Mattel, Monte Cook Games, and Peloton) account for approximately 90\% of the global revenue, as of 2023.
This is key to their success in adopting AI-powered products and reducing labor costs.
It is thus no surprise that these corporations have the resources to invest in fast and widespread adoption of LMs.

The evidence for these corporations adopting LMs is also covered by lawsuits and law firms' representations~\citep{mishcon_generative_ai_ip_tracker} and research studies~\citep{mishcon_generative_ai_ip_tracker}.
But, because of their resources and importance in their industries, there are often no lawsuits or research studying the impact of AI in these industries.
For example, we found no lawsuits of research studying the impact of AI on the video or board games and software industries.
In fact, a search using the terms \emph{Artificial Intelligence} and \emph{detection} or \emph{classifier} yields no relevant cases and articles for the video and board games industries, and only one for the software industry.\footnote{\citet{fowler2023falsepositives}.}

\citet{mishcon_generative_ai_ip_tracker} also reveal that the intellectual property practices of these large corporations have a high impact on the fast and widespread adoptions of LMs.
These corporations thus have the resources to protect their AI solutions and are aware of the potential \emph{intellectual property rights} that belong to them.
They are also the main drivers and owners of large datasets, collected and used for training and fine-tuning LMs~\citep{longpre2024large}.
Other works also finds that the owners or creators of models are often not mentioned or considered when training models.
Additionally, the owners or creators of models are often not considered even when the data used for training or fine-tuning is appropriately credited.
These works also finds that the owners or creators of models are often not considered or are not properly represented when using an API.

\subsection{Economic Impact: Quantifying the Trade-offs}
\label{EconEvidence}

We present evidence from a variety of fields in support of the following claims:

    (i) LMs have increased output and consumer welfare in the short run (though larger among high-income and -educated individuals)---the \emph{Lever of Riches}.
    (ii) The productivity gains have not translated into higher wages, particularly for low-skill, knowledge-intensive workers---the \emph{Gerschenk effect}.
    (iii) Continuous LM usage and integration may lead to long-term productivity gains and higher output per worker---the \emph{Lever of Riches}.
    (iv) Greater output and wealth inequality might emerge over time (at least in the short run), with access to education and labor protection decreasing in turn---the \emph{prisoner's dilemma}.


\subsubsection{The \emph{Lever of Riches}}

Seeing AI as a \emph{Lever of Riches} also suggests that AI's economic impact might be more beneficial than catastrophic.
There are works which argue that LMs can boost global output, with \citet{boiko2023emergent} highlighting LMs' potential to generate scientific and artistic breakthroughs.
The use of LMs can increase output by:    \emph{A new idea may be created}~\citep{boiko2023emergent} . . . \emph{If this weren't the case, then work would still be done by humans .}
\paragraph{Output.}
The use of LMs increases output by extending access to education and increasing labor utilization (the amount of time workers spend productively contributing to the work . . . A second channel . . . is that AI enhances the collective intelligence of humans . \textbf{Over time,} AI enhances the productivity of human users by making information available to them that would have been otherwise unavailable to them . . . \emph{Higher output .}

The \emph{Bartlett effect} shows how the use of LMs increases output.
The use of LMs allows individuals to utilize the ideas, techniques, and data produced by these tools, often on a mass scale.
This enables more productive users, following the \emph{Bartlett effect}.
The ideas, techniques, and data are often produced by other humans, and thus the use of one LM user leads to productivity gains for others.
The \emph{Bartlett effect} can be expected to occur over time, though its pace and magnitude depend on factors such as ease of use and reliability.

The increased output and wealth effects \emph{may} translate into faster economic growth~\citep{restrepo2023automation}.
The duration and magnitude of this translation depend on the short-run and long-run labor market tightness (i.e., unemployment) and the labor market's \emph{complementarity} properties~\citep{restrepo2023automation}.
The latter depend on the labor market's \emph{complementarity} properties.
The labor market can be more or less absorptive to technological advancements . \emph{The} implication is that labor markets, together with technological advances, may exhibit diminishing marginal rates of output gains, even when marginal rates of productivity continue to increase with more advanced technologies.\footnote{Labor markets can be \emph{labor-adaptive} or \emph{labor-complementing}, depending on whether technological improvements reduce or increase labor \emph{utilization rates} (output/workers).}
Some works show that labor utilization rates (output/workers) tend to decrease with labor-market tightness.
This is because, the closer the economy is to its full-employment level, the less room there is to absorb additional labor.
Similarly, empirical evidence from \citet{restrepo2023automation} shows that labor utilization rates have been falling over the past two decades.

\paragraph{Consumers.}
The use of LMs increases consumer welfare by: \emph{This leads the use of a AI model to improve consumer welfare with respect to a model . . . \textbf{In the short run}, the use of AI leads to increased consumer welfare with respect to both customized and off-the-shelf models; . . . .}\footnote{\emph{In the short run} means that LMs benefit consumers in their use cases, without considering their long-term effects on labor, output, or welfare. It also implies that any benefits are immediate and do not account for long-term impacts or sustainable growth.} consumed AI goods. \textbf{In the long run .} \textbf{Over time,} consumer welfare with respect to customized AI solutions . \emph{The use of AI contributes to increased output, which .}

The \emph{Coase theorem} shows how property right allocations can benefit LMs' creators and developers~\citep{coase2013problem}.
\citet{lemley2020fair} argues that LMs and ML algorithms do not require copyright protection to be created, but \citet{lemley2020fair} argues that they \emph{do} require copyright protection to be properly incentivized.
The \emph{Coase theorem} provides an argument for how technological improvements are beneficial to the technology user~\citep{coase2013problem}.
The use of a technology requires different property right allocations that, without intellectual property, may not be practical to realize or offer to technology users.
Intellectual property enables their realization, and the resulting \emph{gain} in consumer welfare \emph{exceeds} the \emph{cost} of \emph{using} the technology (the royalty paid for the IP)\footnote{The \emph{gain} in consumer welfare is the difference between the benefits and the cost of using the technology.}  (i) \emph{Higher education:} Instructors . \emph{Higher, $40\%$ royalty.}? Use LMs in teaching and .  The evidence . In contrast, . Thus . The argument ., with $40\%$ of instructors in Australia, the UK, and the US being either mandated or encouraged to not use AI in essays~\citep{harvard2025chatgpt}. (ii) \emph{Low income:} The evidence . The \emph{rich} are also . The \emph{better educated} also benefit more . The \emph{marginal income} of . The . The \emph{better educated} also benefit more . The \emph{marginal income} of the . The . The \emph{younger} use . The . The \emph{marginal income} of . The . The . The \emph{marginal income} of . The . The . (iii) \emph{Long-term use:} The . The . The \emph{older} use . The . The \emph{marginal income} of . The . The \emph{older} use . The . The \emph{marginal income} of . The . The .


\subsubsection{The \emph{Gerschenk effect}}

The \emph{Gerschenk effect} shows how the use of LMs destroys comparative advantages for human workers in knowledge and skill-intensive jobs, while complementing comparative advantages for human workers in low-skill jobs.
The use of LMs in a field reduces comparative comparative advantages of human workers in that field, especially in knowledge and skill-intensive jobs.
The \emph{Leisure Class theory} emphasizes the importance of the ability to use the ideas and techniques developed by others in reducing the cost of production.

\bibliographystyle{plainnat}
\bibliography{icais2025_conference}


\end{document}
